% theme: similar_triangles_side_ratios
\MajorQuestionLesson{相似な三角形と辺の比}
\SetRowSpace{3mm}

\Needspace{80mm}
\TypeQuestion{}{similar_triangle_ratio_nested}
\FigProblemBlock{53mm}{fig/lesson_03_nested}{
\Text{図で、点 $\mathrm{D},\mathrm{E}$ はそれぞれ辺 $\Lseg{AB},\Lseg{AC}$ 上にあり、$\Lseg{DE}\Parallel\Lseg{BC}$ である。長さの単位は $\mathrm{cm}$ とする。}
\QQ{相似な三角形を、対応する頂点の順に書きなさい。また、相似条件を書きなさい。}{}
\QQ{$\Lseg{AD}:\Lseg{AB}$ を求めなさい。}{}
\QQ{線分 $\Lseg{EC}$ の長さを求めなさい。}{}
}{
\QQ{線分 $\Lseg{DE}$ の長さを求めなさい。}{}
}

\Needspace{80mm}
\TypeQuestion{}{similar_triangle_ratio_crossing}
\FigProblemBlock{53mm}{fig/lesson_03_crossing}{
\Text{図で、$\Lseg{AB}\Parallel\Lseg{CD}$ であり、線分 $\Lseg{AC},\Lseg{BD}$ は点 $\mathrm{O}$ で交わっている。長さの単位は $\mathrm{cm}$ とする。}
\QQ{相似な三角形を、対応する頂点の順に書きなさい。また、相似条件を書きなさい。}{}
\QQ{$\Tri{OAB}$ と $\Tri{OCD}$ の相似比を求めなさい。}{}
\QQ{線分 $\Lseg{OD}$ の長さを求めなさい。}{}
}{
\QQ{線分 $\Lseg{CD}$ の長さを求めなさい。}{}
}

\Needspace{80mm}
\TypeQuestion{}{similar_triangle_ratio_shared_angle}
\FigProblemBlock{53mm}{fig/lesson_03_shared_angle}{
\Text{図で、点 $\mathrm{D}$ は辺 $\Lseg{BC}$ 上にあり、$\Ang{BAD}=\Ang{ACB}$ である。長さの単位は $\mathrm{cm}$ とする。}
\QQ{相似な三角形を、対応する頂点の順に書きなさい。また、相似条件を書きなさい。}{}
\QQ{$\Tri{ABD}$ と $\Tri{CBA}$ の相似比を求めなさい。}{}
\QQ{辺 $\Lseg{BC}$ の長さを求めなさい。}{}
}{
\QQ{線分 $\Lseg{AD}$ の長さを求めなさい。}{}
}

\LessonQuestionBreak

\Needspace{80mm}
\TypeQuestion{}{similar_triangle_ratio_find_length}
\Text{次の図で、$x$ の値を求めなさい。同じ印をつけた角は等しく、同じ平行の印をつけた線分は平行である。長さの単位は $\mathrm{cm}$ とする。}
\QQRowThree{\raisebox{-\height}{\LMFigInline[\linewidth]{fig/lesson_03_ratio_01}}}{}{\raisebox{-\height}{\LMFigInline[\linewidth]{fig/lesson_03_ratio_02}}}{}{\raisebox{-\height}{\LMFigInline[\linewidth]{fig/lesson_03_ratio_03}}}{}
\QQRowThree{\raisebox{-\height}{\LMFigInline[\linewidth]{fig/lesson_03_ratio_04}}}{}{\raisebox{-\height}{\LMFigInline[\linewidth]{fig/lesson_03_ratio_05}}}{}{\raisebox{-\height}{\LMFigInline[\linewidth]{fig/lesson_03_ratio_06}}}{}
\QQRowThree{\raisebox{-\height}{\LMFigInline[\linewidth]{fig/lesson_03_ratio_07}}}{}{\raisebox{-\height}{\LMFigInline[\linewidth]{fig/lesson_03_ratio_08}}}{}{\raisebox{-\height}{\LMFigInline[\linewidth]{fig/lesson_03_ratio_09}}}{}
\QQRowThree{\raisebox{-\height}{\LMFigInline[\linewidth]{fig/lesson_03_ratio_10}}}{}{\raisebox{-\height}{\LMFigInline[\linewidth]{fig/lesson_03_ratio_11}}}{}{\raisebox{-\height}{\LMFigInline[\linewidth]{fig/lesson_03_ratio_12}}}{}
\QQRowThree{\raisebox{-\height}{\LMFigInline[\linewidth]{fig/lesson_03_ratio_13}}}{}{\raisebox{-\height}{\LMFigInline[\linewidth]{fig/lesson_03_ratio_14}}}{}{\raisebox{-\height}{\LMFigInline[\linewidth]{fig/lesson_03_ratio_15}}}{}
